\section{The character condition throughout the middle range}

% BEGIN M81B1
\paragraph{M81B1. The ambient jump window at every middle-range place.}
Use the original family hub $X$ of F78D2. Fix a twisted index $i$,
$p=\ell_i$ outside the two type primes, in the native range
$\mathrm{RangeAD}(L,X.\mathrm{exc},p)$. Thus $p$ is outside the
host small set, $p\ne3,q,d$, and either $p\le\ell_\star$ or
$p\in X.\mathrm{exc}$. Take any $\pi\in Z_{\rm good}(i)$
and any $w\mid p$. These are the simultaneous guards for the
split and inert branches of $\mathrm{StarAD}$.

The conductor criterion gives $p\nmid N$ and $p\ne2$.
The argument of S79P1 applies to this $w$ regardless of the
declaration: the host embeds in the conductor-$N$ cyclotomic field,
unramified at $p$. Its completed bridge $K$ is therefore unramified,
$K=K_0$, and $d_w=[K:\mathbb Q_p]\le [E:\mathbb Q]=4n$.
The bridge inertia is $I_K$ by transport of the valuation and its
residue action. Every good coefficient extension $L'$ tested in
$\mathrm{admTB}$ is good at level one: $K\subset K_0(\mu_p)$,
and the given splitting of the higher cyclotomic layer implies
the same at level one by restriction of embeddings.

Put $H=\lfloor2b\rfloor$, $T=(t-1)/2$ and $j=j_i$.
At a label $\tau:K\to L'$ use the away-from-type reading of
$\mathrm{AdmJ}(L)$ on $\mathrm{GoodR}(X.r_i)$. Before removing
any summands, subtraction of the $\chi_1$-weight in S79P3 gives
the cloud $T+v+v_a$, the lows $T+v-b_a$ and the string
$0,\ldots,t-1$, with $v=e_{j-1}$ at the archimedean place.
At the conjugate orientation reflect all entries by $z\mapsto2T-z$.
The original dictionary applies over the specified coefficient
extension, as in S80J1--S80J2, so these lists cover the jumps of the actual twisted
ambient representation $R_\pi\chi_1^{-1}$. They contain at most
$(k-1)+m+t=n_V$ distinct integers, counted together.
The parameter bounds give $0\le v\le b+\eta+1$,
$0\le v_a\le S_c$ and $|v-b_a|\le\eta$ (G79E3).
Hence both signs of the cloud and lows have absolute value at most
$b+\eta+1+S_c+T$. The same bound contains the string.
Since $S_c+\eta+t+2\le b$ and $T\le t$, every integral
entry has $|z|+1\le H$ and $t\le H$. For the kernel lattices
tested below take the union of these integer lists over all
$d_w$ labels. The graded filtration decomposes over
$L'\otimes K=\prod_\tau L'$, so a jump of the total filtration
must occur on at least one label. Thus the union bounds the
total jump set as well, with cardinality at most $4nn_V$.
% END M81B1

% BEGIN M81B2
\paragraph{M81B2. The kernel exponent on the actual unramified bridge.}
Put $N=4nn_V$, $P=\lfloor\log_p(2H)\rfloor$. Since $K=K_0$,
the translated cyclotomic polynomials are Eisenstein over $K$.
The good level-one field has degree $p-1$ over $K$, and the
level-$j$ field has degree $p^{j-1}$ over the level-one field
for $j\ge1$. Hence $g_K=0$ and $\mathrm{vdegMs}(j)=j-1$.
For each crystalline $Y$ with its labelled jumps among those in
M81B1, take $M_0=h_0=H$. The level
$m_0=3+\min\{v\ge0:2H\le p^v(p-1)\}$ is at most $4+P$.
For $0\le i'\le e<t\le H$ and each jump $u\ne i'$, the
nonzero integer $u-i'$ has absolute value at most $2H$, so
$v_p(u-i')\le P$. Each inner sum in $\mathrm{cLL}$ costs at
most $N(4+2P)$. The initial degree costs $m_0-1$, and the
$e+1\le t$ terms add two to the maximum of one and their
respective inner sum. The remaining $\mathrm{cTB}$ corrections
are $e+2+v_p(e)+\lfloor e/(p-1)\rfloor+1+1+v_p(e+1)$,
at most $4t+3$ with the native $v_p(0)=0$. Consequently
\[
 \mathrm{cTBY}(Y,H,H,e)\le
 (3+P)+t\{3+N(4+2P)\}+4t+3=X.c_B(i).
\]
The equality is the defining $\mathrm{cBof}$ of DFD8 and C78B3
on this hub. Together with M81B1 this gives
$\mathrm{JointConditionTameBound}$ at its stored exponent.

For the first clause of $\mathrm{admTB}$ take the tested
crystalline representation $\widetilde Y$, fixed quotient
functional and integral lift $y$ of its generator. The
kernel has lattice $T=\ker(\pi_0)\cap\widetilde T$,
where $\pi_0$ denotes that functional. Its quotient lattice is
the trivial integral line because $\pi_0(y)=1$ and the functional
is integral on $\widetilde T$. Strict crystalline subobject
comparison (H80M1) puts the kernel's jumps among the tested
ambient jumps and makes the kernel crystalline. W79H1 and W79H5
give crystallinity over the good level-one subgroup and the
required $\mathrm{HypW}$ at heights $H,H$ on the chosen
comparison. The actual integral connecting cocycle is $gy-y$
viewed in $T$. W81T6 and K81C6 applied to it show that its
image under every tested equivariant map to the inverse
cyclotomic divisible line becomes finite after multiplication
by $p^{\mathrm{cTBY}(Y,H,H,e)}$. Multiplication by the remaining
power gives the first clause at $X.c_B(i)$.
% END M81B2

% BEGIN M81B3
\paragraph{M81B3. The possible cyclic obstruction and all remaining character clauses.}
Write $O=\mathcal O_{L'}$, $V=L'(-e)$, let $T_e$ be its given
stable lattice, and write $A_e=V/T_e$. Every allowed realization
is of this form by an integral choice of basis of the line.
For every $e\ge0$, Tate duality identifies $H^2(K,T_e)$ with
the Pontryagin dual of the invariants of the divisible twist
of weight $e+1$. Here the trace-dual coefficient lattice contributes
the inverse different, a free rank-one fractional $O$-ideal.
The cyclotomic image contains $1+p$. Its $(e+1)$st power differs
from one, so the invariants lie in the kernel of a nonzero scalar
on one copy of $L'/O$. This kernel and all its submodules are
finite cyclic over the DVR. The Pontryagin dual of $O/I$ is
also cyclic: using the trace-dual lattice $O^\vee$, it is
$I^{-1}O^\vee/O^\vee$ for $I\ne0$. This includes the zero
module when $I=O$. Thus the possible integral $H^2$ contributes
at most one cyclic summand. Rational $H^2(K,V)$ vanishes by
duality since the positive dual twist has zero rational invariants.

The local Euler formula gives
$\dim_{L'}H^1(K,V)=d_w+\mathbf1_{e=0}$.
For these crystalline lines the Bloch--Kato dimension is
$\mathbf1_{e=0}$: their de Rham filtration has full degree-zero
step, and only $e=0$ has a fixed vector.
For $e=0$ the unramified homomorphisms give the entire rational
finite line, since their Witt Frobenius primitive is crystalline
(K81C6). Every unramified homomorphism to $A_0$ lifts rationally
by lifting its Frobenius value and composing with the projection
of the residue Galois group to $\mathbb Z_p$.
Thus finite torsion cocycles for $e=0$ are precisely unramified
ones up to coboundaries. For $e>0$ they are precisely coboundaries,
since the rational finite space is zero. These are the two
native cocycle assertions, using the inertia identification of M81B1.

The image of $H^1(K,T_e)$ in $H^1(K,V)$ is a full $O$-lattice.
Indeed it spans since a continuous rational cocycle has bounded
image on the compact group. It is compact: the integral
cohomology is the inverse limit of the cohomologies of the finite
lattice quotients, which are finite by local finiteness, and the
degree-zero inverse system also consists of finite groups, so
there is no derived-limit term. For an $O$-uniformizer $\lambda$,
the quotient $H^1(K,T_e)/\lambda H^1(K,T_e)$ injects into the
finite group $H^1(K,T_e/\lambda T_e)$. Lift finitely many
generators of that quotient. Successive
approximation with these lifts converges in the compact inverse
limit, since powers of the uniformizer kill each fixed finite
lattice quotient. This gives finite generation over $O$.
Its projected image in
$H^1(K,V)/H^1_f(K,V)$ is likewise a full lattice.
The long exact sequence for $T_e\subset V$, followed by dividing
by the image of the rational finite part, consequently gives
\[
 0\longrightarrow (L'/O)^{d_w}\longrightarrow
 H^1(K,A_e)/H_f^1(K,A_e)\longrightarrow H^2(K,T_e)
 \longrightarrow 0 .
\]
Indeed the middle map before quotient has kernel the rational
cohomology modulo the integral lattice and cokernel zero by
rational $H^2=0$. Taking $p^a$-torsion is left exact. Its length
in the left term is $d_wa\operatorname{length}_O(O/p)$,
and in the right term it is at most
$a\operatorname{length}_O(O/p)$ by cyclicity.
The same bound $(d_w+1)a\operatorname{length}_O(O/p)$ therefore
holds in the middle. At $a=X.c_B(i)$, the quotient of cocycles
whose multiple is finite by the finite cocycles injects in
this torsion group, since all coboundaries are already finite.
This gives the native length clause even when $p-1$ divides $e+1$.

Together M81B1--M81B3 prove $\mathrm{admTB}$ at every tested
middle-range point on this hub, for both the split and the inert
case, also on restriction to the rows of a split declaration.
The rank-representation row estimates and the split-column
conditions in $\mathrm{StepAD}$ require their own arguments.
% END M81B3
